\section{System Overview}

The repository is organized as a reproducible MATLAB pipeline. The top-level script \texttt{main.m} loads configuration, resets generated outputs, generates synthetic datasets, runs all filter configurations, and exports tables and figures. The modular scripts \texttt{generate\_dataset.m}, \texttt{run\_filters.m}, and \texttt{generate\_outputs.m} separate dataset creation, filtering, and post-processing.

\subsection{Pipeline}

For each configured scenario, \path{generate_scenario_dataset.m} creates a wall-segment map, beacon definitions, a wall-validating ground-truth trajectory, noisy prediction inputs, beacon measurements, LiDAR scans, and synthetic LiDAR pose measurements. The data is then loaded by \texttt{run\_filter\_case.m}, which evaluates five cases using the same truth and measurement files. Finally, \texttt{collect\_results\_tables.m} computes RMSE and sensor availability metrics, and \texttt{generate\_project\_figures.m} exports the result figures used in this manuscript.

\subsection{Ground Robot State}

The ground robot state is a 2D NavState matrix with heading, position, and velocity. The stored truth arrays use \texttt{truth.p} and \texttt{truth.v} as \(2 \times N\) matrices, \texttt{truth.theta} as a length-\(N\) heading sequence, and \texttt{truth.X} as \(4 \times 4 \times N\) NavState matrices. No accelerometer bias, gyroscope bias, clock state, or beacon state is estimated by the EKF.

\subsection{Environment Representation}

The environment is a known 2D floor plan represented by line-segment walls. Each wall stores endpoints, thickness, material factor, and a text label. Openings are represented by gaps between wall segments. The simulator uses wall geometry for trajectory validation, beacon attenuation, and LiDAR ray casting. The filter does not receive wall-crossing counts or hidden attenuation values.

\subsection{Beacon and LiDAR Simulation}

Beacon simulation is implemented in \texttt{simulate\_beacon\_measurements.m}. Each beacon is fixed in the map and has a range limit, penetration budget, RSSI reference, RSSI threshold, and range-noise parameters. The simulator rejects measurements when range, attenuation, RSSI, or random dropout tests fail. The filter receives only timestamp, beacon ID, measured range, RSSI, and availability.

LiDAR simulation is implemented in \texttt{simulate\_lidar\_measurements.m}. The scan model ray-casts 181 beams across a \(270^\circ\) field of view with a 12 m maximum range. The EKF does not update on raw beams. Instead, the simulator produces a synthetic pose measurement and covariance whose standard deviations depend on scan quality. 

\subsection{Estimation Backend}

The estimation backend is an error-state EKF. Prediction is driven by noisy body-frame acceleration and yaw-rate measurements. Corrections are applied using scalar beacon range residuals and three-dimensional LiDAR pose residuals. The covariance is maintained in the local tangent ordering
\begin{equation}
    \delta x =
    \begin{bmatrix}
    \delta \theta & \delta p_x & \delta p_y & \delta v_x & \delta v_y
    \end{bmatrix}^{T}.
\end{equation}
After each correction, \texttt{ekf\_tangent\_update.m} maps the translational parts of the correction into the local group frame and composes the result onto the nominal NavState using \texttt{navstate\_exp.m}.
